% page 150 of the facsimile (image p-149 of the scan)
% block 157.5 x 246.3 mm ; left margin 8.0 mm
\pgc{-12.700mm}{0.000mm}{
\l{\cel{6}140}
\l{}
\l{}
\l{\sou{epizootio}. Morbo epidemiala, o kontagiala, qua afektas tota}
\l{\cel{3}klaso de animali domestika en un foyo. - DEFIRS.}
\l{}
\l{\sou{epodo}.- I. En la poezio Greka e Latina : la dividuro triesma}
\l{\cel{3}e lasta di la odo, qua venas pos la strofo e la anti-strofo}
\l{\cel{3}- II. Verso duesma (iambo quar-silaba) di distiko di qua}
\l{\cel{3}la verso unesma esas iambo sis-silaba. - III. Mikra poemo}
\l{\cel{3}ek distiki tala. - DEFIRS.}
\l{}
\l{\sou{epoko}. Tempo-parto qua havas kom distingivo ulo historio-}
\l{\cel{3}importoza. - DEFIRS.}
\l{}
\l{\sou{epoleto}. Insigno armeanala, an la shultro, ek bendo di qua}
\l{\cel{3}la extremajo mi-cirkla garnisesas ye franjo. - DEFIRS.}
\l{}
\l{\sou{epruveto}. Tubo ek vitro o kristalo, di qua un extremajo esas}
\l{\cel{3}klozita, qua uzesas por mezurar la volumino, la elastik-}
\l{\cel{3}eso di la gasi, por agar experimenti kemiala diversa, por}
\l{\cel{3}determinar la koncentreso-grado di ula liquidi, e c. -FIS.}
\l{}
\l{\sou{epuliso}. (\sou{patol.)}\cel{2}Surkreskajo di la jenjivo.}
\l{}
\l{\sou{equaciono.}\cel{1}(\sou{algebro}).Formulo qua expresas egaleso inter du}
\l{\cel{3}quanti algebrala, e qua kontenas ordinare un o plura ne-}
\l{\cel{3}konocaji. - EFIS.}
\l{}
\l{\sou{equatoro}. I. Granda cirklo di la Ter-sfero, perpendikla ye}
\l{\cel{3}la rotac-axo, e qua dividas la globo a du mi-sferi. -II.}
\l{\cel{3}Granda cirklo diSuno e di la planeti, perpendikla ye lia}
\l{\cel{3}rotac-axo. - III. Granda cirklo qua rezultas de la inter-}
\l{\cel{3}seko di la plano di la Ter-equatoro e di la cielo-sfero.}
\l{\cel{3}- DEFIRS.}
\l{}
\l{\sou{equi-}\cel{2}Prefixo uzata exkluzive en la vorti ciencala, quasigni-}
\l{\cel{3}fikas \textquotedbl{}egala\textquotedbl{}.}
\l{}
\l{\sou{equilibrar}. (\sou{netrans.}) I. Esar en la situeso di korpo, di}
\l{\cel{3}punto materiala, qua - atraktate da forci inter-egala ed}
\l{\cel{3}inter-kontrea - permanas sen movo. - II. Esar egale repar-}
\l{\cel{3}tisita (pri fluido en ta o ca medio). - II. (\sou{pri la kozi}).}
\l{\cel{3}Esar distributita saje. - EFIRS.}
\l{}
\l{\sou{equinoxo}. Singla ek la du periodi di la yaro dum qua (pro ke}
\l{\cel{3}Suno pasas en la equatoro), la dio-duro egalesas la nokto-}
\l{\cel{3}duro sur la tota Ter-globo. - DEFIS.}
\l{}
\l{\sou{equipar}. (trans.) I. Provizar (navo) per omno necesa por la}
\l{\cel{3}manovro, e por la nutro e la armizo di la homi embarkota.}
\l{\cel{3}- II. Provizar (armeani) per la kozi necesa, vesti, armi,}
\l{\cel{3}e c. - III. Provizar (mashino, aparato) per omna kozi}
\l{\cel{3}necesa por lua funciono. - DEFIR.}
\l{}
\l{\sou{equi}polar. (\sou{trans.}) (\sou{matem.}) Equilibrigar segun la metodo}
\l{\cel{3}da Giusto Bellavitis.}
}
